\section{Preliminaries}\label{Section:Preliminaries}

\subsubsection{Notation.}
Let $\mathbb K$ be a field with algebraic closure
$\overline{\mathbb K}$.
Throughout, $\mathcal{F} = \{f_1, \ldots, f_n\} \subset \mathbb{K}[\avec][\xvec]$ is a system of $n$ equations in $n-1+k$ variables, viewed as polynomials in the eliminated variables $\xvec=(x_1,\ldots,x_{n-1})$ with coefficients in $\mathbb{K}[\avec]$, where $\avec=(a_1,\ldots,a_k)$ are the retained parameters.
The total degree of $f$ is denoted by $\deg(f)$ and its degree in $x_i$ by $\deg_{x_i}(f)$; $\langle \mathcal{G} \rangle$ denotes the ideal generated by a set $\mathcal{G}$ of polynomials.
All notation is summarized in Table~\ref{tab:notation}.

\begin{table}[t]
	\centering
	\small
	\caption{Notation used throughout the paper.}
	\label{tab:notation}
	\setlength{\tabcolsep}{4pt}
	\renewcommand{\arraystretch}{1.05}
	\begin{tabular}{@{}l@{\quad}p{0.8\textwidth}@{}}
		\toprule
		\textbf{Symbol} & \textbf{Meaning} \\
		\midrule
        $\mathcal{F}$ & Polynomial System $\mathcal{F}=\{f_1, \dots, f_n\} \subset \mathbb{K}[\avec][\xvec]$ \\
		$\xvec$, $\overline{\xvec}$, $\avec$ & $n-1$ eliminated variables, their dual variables, $k$ retained parameters \\
        $X$, $\overline X$ & row vectors of monomials in $\xvec$ and $\overline{\xvec}$, respectively \\
		$C$, $\widehat{C}$ & cancellation matrix~\eqref{eq:Cancellation Matrix} and its refined form~\eqref{eq:Dixon_polynomial_improve} \\
		$\Delta_{\mathcal{F}}$, $\mathcal{M}$ & Dixon polynomial~\eqref{eq:Dixon_polynomial} and Dixon matrix (\hyperref[dixon_step2]{S2}), $\Delta_{\mathcal{F}}=X\mathcal{M}\overline{X}^{T}$ \\
		$D(\mathcal{F})$ & Dixon matrix size of $\mathcal{F}$: the maximum of its row and column counts \\
		$\mathcal{M}'$, $\rho$ & maximal-rank square submatrix of $\mathcal{M}$, of size $\rho=\operatorname{rank}\mathcal{M}\le D(\mathcal{F})$ \\
		$\mathcal{I}$, $J$, $d_{\mathcal{I}}$ & $\langle\mathcal{F}\rangle$; elimination ideal $\mathcal{I}\cap\mathbb{K}[\avec]$; degree of $\mathcal{I}$ \\
		$I_{d_1,\ldots,d_n}$ & Iterated-Simplex polynomial~\eqref{eq:Iterated_Simplex_Polynomial} \\
        $C_n^{(d)}$ & Fuss--Catalan number~\eqref{eq:fuss_catalan} \\
        $\mathcal X$, $|\mathcal X|$ &
        variable block and its cardinality (\hyperref[subsec:method2]{M2}) \\
		$\omega$ & matrix multiplication exponent, $2\le\omega<3$ \\
		\bottomrule
	\end{tabular}
\end{table}

\subsection{Resultants}

Resultants are classical tools in elimination theory used to determine whether a polynomial system has a common solution. We use the projection-operators viewpoint of~\cite{DBLP:conf/issac/KapurS97}.

\begin{definition}[Resultant~\cite{DBLP:conf/issac/KapurS97}]\label{Definition:Resultant}
	Let $\xvec=(x_1, \dots, x_{n-1})$ be the variables to eliminate and $\avec=(a_1, \dots, a_k)$ the retained parameters. For a system $\mathcal{F}=\{f_1, \dots, f_n\} \subset \mathbb{K}[\avec][\xvec]$, let $\mathcal{I} = \langle \mathcal{F} \rangle \subset \mathbb{K}[\avec, \xvec]$ and $J := \mathcal{I} \cap \mathbb{K}[\avec]$. Polynomials in $J$ are called \emph{projection operators} and vanish on the parameter values for which $\mathcal{F}$ has a common solution over $\overline{\mathbb{K}}$. If $J$ is principal, its generator (up to scalar) is the \emph{resultant} of $\mathcal{F}$.
\end{definition}

The Sylvester resultant~\cite{cox2005using} for two univariate polynomials of degrees $d_1,d_2$ produces a matrix of size $d_1+d_2$; the B\'ezout--Cayley resultant~\cite{painvin1874note,cox2005using} reduces this to $\max(d_1,d_2)$. The Macaulay resultant~\cite{macaulay1902some} generalizes this to several multivariate polynomials but can produce very large matrices. The Dixon construction~\cite{Dixon1909} is a multivariate extension of the B\'ezout-style approach: it preserves the compactness of the B\'ezout--Cayley matrix while extending it from two to $n$ polynomials, often yielding much smaller matrices than Macaulay.

\begin{remark}\label{rem:projective_vs_affine}
	For $k > 1$, $J$ need not be principal, and determinant-based methods compute a specific nonzero element $R \in J$. The affine determinantal construction is defined as the determinant of a maximal-rank submatrix of the Dixon matrix $\mathcal{M}$ of Section~\ref{dixon_step2} and vanishes on the projection of the affine solution set, but may include extraneous factors not essential for the projection. In this paper, we refer to this determinant $R$ as the resultant in a slightly informal sense, rather than as a projection polynomial.
\end{remark}

\subsection{Dixon Resultant}\label{sec2.3:Dixon}
The Dixon resultant method~\cite{Dixon1909,DBLP:conf/issac/KapurSY94} eliminates several variables simultaneously via a determinant. The procedure has four steps~\cite{DBLP:conf/issac/KapurSY94,JinaduMonagan2025}.

\subsubsection{Step 1: Compute the Dixon polynomial.}\label{dixon_step1}
Introduce $n-1$ auxiliary variables $\overline{x}_1, \ldots, \overline{x}_{n-1}$ and define the $n \times n$ cancellation matrix
\begin{equation}
	C(\xvec, \overline{\xvec}) = 
	\begin{bmatrix}
		f_1(x_1, x_2, \dots, x_{n-1}) & \cdots & f_n(x_1, x_2, \dots, x_{n-1}) \\
		f_1(\overline{x}_1, x_2, \dots, x_{n-1}) & \cdots & f_n(\overline{x}_1, x_2, \dots, x_{n-1}) \\
		f_1(\overline{x}_1, \overline{x}_2, \dots, x_{n-1}) & \cdots & f_n(\overline{x}_1, \overline{x}_2, \dots, x_{n-1}) \\
		\vdots & \ddots & \vdots \\
		f_1(\overline{x}_1, \overline{x}_2, \dots, \overline{x}_{n-1}) & \cdots & f_n(\overline{x}_1, \overline{x}_2, \dots, \overline{x}_{n-1})
	\end{bmatrix}.
	\label{eq:Cancellation Matrix}
\end{equation}
The Dixon polynomial is
\begin{equation}
	\Delta_{\mathcal F}(\xvec,\overline{\xvec}) = \frac{\det(C(\xvec, \overline{\xvec}))}{\prod_{i=1}^{n-1} (x_i - \overline{x}_i)}.
	\label{eq:Dixon_polynomial}
\end{equation}

Instead of performing an explicit symbolic division, we use the refined construction from~\cite{DBLP:conf/casc/JinaduM22}
\begin{equation}
	\mathrm{Row}(\widehat{C}_1) = \mathrm{Row}(C_1), \qquad
	\mathrm{Row}(\widehat{C}_{i+1}) = \frac{\mathrm{Row}(C_{i+1}) - \mathrm{Row}(C_i)}{x_i - \overline{x}_i},
	\label{eq:Dixon_polynomial_improve}
\end{equation}
for $i=1,\ldots,n-1$. Each numerator is divisible by $x_i-\overline{x}_i$ in $\mathbb{K}[\xvec,\overline{\xvec},\avec]$, so the quotient remains polynomial; the Dixon polynomial is $\Delta_{\mathcal F}(\xvec,\overline{\xvec})=\det(\widehat{C})$. This row-wise difference-quotient construction is used throughout the paper.

\subsubsection{Step 2: Construct the Dixon matrix.}\label{dixon_step2}
The Dixon polynomial admits a bilinear form $\Delta_{\mathcal F}(\xvec,\overline{\xvec}) = X \mathcal{M} \overline{X}^{T}$, where $X$ and $\overline{X}$ list monomials in $\xvec$ and $\overline{\xvec}$, respectively. The coefficient matrix $\mathcal{M}$ is the \textbf{Dixon matrix}: its rows are indexed by the $\xvec$-monomials and its columns by the $\overline{\xvec}$-monomials occurring in $\Delta_{\mathcal F}$, so it is in general rectangular and rank-deficient. Alternatively, $\mathcal{M}$ can be constructed directly via the recursive block method of Zhao and Fu~\cite{Zhao2005extfast}, particularly relevant when few variables are eliminated and degrees are high~\cite{DBLP:journals/ijcm/QinWTJ17}.

\subsubsection{Step 3: Extract a maximal-rank submatrix.}\label{dixon_step3}
In practice, the Dixon matrix $\mathcal{M}$ is often rectangular or rank-deficient, so its full determinant is unavailable or identically zero. Following the extension of Dixon elimination to singular matrices by KSY~\cite{DBLP:conf/issac/KapurSY94}, we extract a maximal-rank square submatrix $\mathcal{M}'$ of size $\operatorname{rank}\mathcal{M}$ whose determinant is nonzero. Following~\cite{DBLP:conf/casc/JinaduM22,JinaduMonagan2025}, such a submatrix is typically identified efficiently after evaluating $\mathcal{M}$ at a generic point of $\mathbb{K}$.

\begin{remark}
The validity of using $\det(\mathcal M')$ relies on the
KSY condition~\cite{DBLP:conf/issac/KapurSY94}, also called the RSC criterion~\cite{Albrecht2008}: there exists a column whose monomial label is nonzero on the solution set under consideration and whose deletion
lowers the rank over $\mathbb K(\avec)$. Failures of this condition
appear rare in practice~\cite{DBLP:conf/issac/KapurSY94},
and Albrecht~\cite[Section~4.2.3]{Albrecht2008} addresses
the nontermination issue in the RSC procedure through
bounded retries and recursive specialization.
\end{remark}

\subsubsection{Step 4: Compute the Dixon resultant.}\label{dixon_step4}
Once $\mathcal{M}'$ is obtained, $\det(\mathcal{M}')$ is the Dixon resultant in $\mathbb{K}[\avec]$.
Under the KSY condition, it remains a valid projection polynomial for a singular Dixon matrix: it vanishes whenever the original system has a common solution in the eliminated variables. In general, $\det(\mathcal{M}')$ may contain extraneous factors.


\subsubsection{Size of the Dixon matrix.}
The size of the Dixon matrix is a key parameter in the complexity analysis. A classical variable-wise upper bound~\cite{DBLP:conf/issac/KapurSY94,DBLP:journals/ijcm/QinWTJ17} is
\begin{equation}
	B_{var} = (n-1)! \prod_{i=1}^{n-1} m_i,
	\label{eq:Dixon_var_bound}
\end{equation}
where $m_i = \max_{1\leq j\leq n}\deg_{x_i}(f_j)$ is the maximum degree in $x_i$ over all polynomials in $\mathcal{F}$. Kapur and Saxena~\cite[Theorem~4.3]{DBLP:conf/stoc/KapurS96} give a volume-based asymptotic estimate for multi-homogeneous systems. In the equal-total-degree setting, their volume expression is $(dn)^{n-1}/n!$. We use its integer part as the comparison bound
\begin{equation}
	B_{MH} = \left\lfloor \frac{(dn)^{n-1}}{n!} \right\rfloor.
	\label{eq:Dixon_volume_bound}
\end{equation}
Section~\ref{Section: Complexity} derives exact generic support formulas for the total-degree setting and independently verifies the explicit inequality $D(\mathcal{F})\le B_{MH}$ from the Fuss--Catalan formula.

\subsubsection{Influence of extraneous factors.}
The determinantal output may contain factors unrelated to the affine solution set.

\begin{definition}[Extraneous factor]\label{def:extra_factor}
	Let $Z$ be the Zariski closure of the projection of $V(\mathcal{I})$ onto the retained parameters. An irreducible factor $e$ of the determinantal elimination polynomial is an \textbf{extraneous factor} if $V(e)\not\subseteq Z$.
\end{definition}

In the absence of such factors and of excess multiplicities, the total degree of the eliminant is bounded by $d_{\mathcal{I}}$, both when $k=1$ and when $k\ge 2$; in either case it is bounded by the B\'ezout bound:

\begin{theorem}[B\'ezout bound~\cite{Heintz1983,Fulton1998}]\label{thm:bezout_bound}
	Let $\mathcal{I} = \langle f_1, \ldots, f_m \rangle \subset \mathbb{K}[x_1, \ldots, x_N]$ with $\deg(f_i)=d_i$, and let $d_{\mathcal{I}}$ denote the degree of $\mathcal{I}$, i.e.\ the number of solutions counted with multiplicity when $\mathcal{I}$ is zero-dimensional, and the degree of the affine variety $V(\mathcal{I})$ in general. Then
	\[
	d_{\mathcal{I}} \leq \prod_{i=1}^m d_i.
	\]
\end{theorem}

In the presence of extraneous factors, the degree of the eliminant may exceed the B\'ezout bound of the underlying system, both inflating computational costs and introducing spurious components unrelated to the actual solution set.

\subsection{Polynomial Matrix Determinant Computation}\label{subsection:poly_mat_det}
We summarize the determinant complexity models used in this paper; full details are in Appendix~\ref{appendix:determinant_methods}.

\subsubsection{Univariate determinants.}
For $P(t)\in\mathbb{K}[t]^{m\times m}$, the Hermite normal form (HNF) method of~\cite{DBLP:journals/jc/LabahnNZ17} has deterministic cost $T^{\mathrm{HNF}}=\widetilde{O}(m^{\omega}\lceil \delta\rceil)$, where $\delta$ is the average column degree.

\subsubsection{Multivariate determinants.}
For $P(\yvec)\in \mathbb{K}[y_1,\ldots,y_v]^{m\times m}$, no single method is universally best: different algorithms can lead to substantially different complexity estimates depending on the matrix size, the number of variables, and the sparsity structure of the entries. We therefore consider several alternatives. Besides direct symbolic expansion by minors~\cite{HorowitzSahni1975,GentlemanJohnson1976} and fraction-free Bareiss elimination~\cite{Bareiss1968}, one may reduce to the univariate case via Kronecker substitution~\cite{Li2007symbolic} and apply HNF, apply evaluation-interpolation on a tensor grid~\cite{HorowitzSahni1975}, or use sparse interpolation~\cite{DBLP:journals/jsc/Huang23}. The corresponding time and space costs are summarized in Table~\ref{tab:determinant_methods}.

In the table notation, $b_i$ bounds $\deg_{y_i}\det(P)$ and fixes the Kronecker substitution base (chosen to avoid overflow in the encoded determinant); $B:=\max_i b_i$; $N:=\prod_{i=1}^{v}(b_i+1)$ is the tensor-grid size used by interpolation; $S$ bounds the number of nonzero terms of $\det(P)$; $U$ bounds the term count of intermediate minors; $L$ is the cost of one determinant probe; and $\mu\le \binom{v+d}{d}$ bounds the dense monomial count of an entry of total degree at most $d$. All complexities are reported in field operations over $\mathbb{F}_q$; derivations are given in Appendix~\ref{appendix:determinant_methods}.

\begin{table}[t]
	\centering
	\small
	\caption{Determinant computation models used in this paper. All complexities are in field operations over $\mathbb{F}_q$.}
	\label{tab:determinant_methods}
	\setlength{\tabcolsep}{4pt}
	\renewcommand{\arraystretch}{1.05}
	\begin{tabular}{p{0.35\textwidth}p{0.29\textwidth}p{0.24\textwidth}}
		\toprule
		\textbf{Method} & \textbf{Time} & \textbf{Space} \\
		\midrule
		Expansion~\cite{GentlemanJohnson1976,HorowitzSahni1975} & $\widetilde{O}(m2^m \mu U)$ & $\widetilde{O}(\binom{m}{m/2}U)$ \\
		Bareiss~\cite{Bareiss1968,GentlemanJohnson1976} & $\widetilde{O}(m^3U^2)$ & $\widetilde{O}(m^2U)$ \\
		Kronecker~\cite{Li2007symbolic} & $\widetilde{O}(m^{\omega}\delta)$ & $\widetilde{O}(m^2\delta)$ \\
		Interpolation~\cite{HorowitzSahni1975} & $\widetilde{O}\!\left(N(L+vB)\right)$ & $\widetilde{O}(N)$ \\
		Sparse~\cite{DBLP:journals/jsc/Huang23}\footnotemark & $\widetilde{O}(LS + S\log q)$ & $\widetilde{O}(S)$ \\
		\bottomrule
	\end{tabular}
\end{table}
\footnotetext{The sparse model also requires $\mathrm{char}(\mathbb{F}_q)\ge B$ (Huang~\cite[Theorem~1]{DBLP:journals/jsc/Huang23}); it therefore applies to large-prime fields but not to extension fields of the form $\mathbb{F}_{2^\ell}$.}

\subsection{Security of AO Primitives}
The sponge construction~\cite{DBLP:conf/eurocrypt/BertoniDPA08} builds hash functions from permutations. The relevant task for algebraic attacks is the constrained-input constrained-output~\cite{bertoni2007sponge} (CICO) problem.

\begin{definition}[CICO problem]
	For a permutation $F\colon \Fq^t \to \Fq^t$, the CICO-$k$ problem is to find $x \in \Fq^{t-k} \times \{0\}^{k}$ and $y \in \Fq^{t-k} \times \{0\}^{k}$ with $F(x) = y$.
\end{definition}

Brute force requires $\sim q^k$ permutation calls; substantially faster algebraic solutions may translate into preimage or collision attacks in the sponge setting.
